% ====================================================================
% DAFTAR PUSTAKA (BIBLIOGRAPHY)
% ====================================================================
% File ini berisi daftar referensi/pustaka yang digunakan dalam penelitian.
% Format: APA Style (sesuai Juknis UNEJ)
%
% PETUNJUK PENGGUNAAN:
% 1. Gunakan file references.bib untuk menyimpan semua referensi
% 2. Uncomment baris di bawah untuk mengaktifkan daftar pustaka
% 3. Compile dengan: pdflatex → bibtex → pdflatex → pdflatex
% ====================================================================

\setlength{\bibsep}{0pt} % Mengurangi spasi antar referensi (opsional)
\renewcommand{\bibname}{%
\begin{center}
    {\normalfont\bfseries\fontsize{14pt}{16pt}\selectfont DAFTAR PUSTAKA}
\end{center}
\vspace*{-1cm}
}
\addcontentsline{toc}{chapter}{DAFTAR PUSTAKA}

% --- MENGGUNAKAN BIBTEX (AKTIF) ---
% Gunakan file daftar_pustaka.bib untuk menyimpan semua referensi
\bibliographystyle{apalike}
\bibliography{daftar_pustaka}

% ====================================================================
% PANDUAN PENGUMPULAN REFERENSI:
%
% 1. MINIMAL REFERENSI:
%    - Juknis: minimal 50% dari 10 tahun terakhir
%    - Umumnya: 20-50 referensi untuk skripsi
%
% 2. JENIS SUMBER YANG BAIK:
%    - Academic journals (50-70%)
%    - Books (10-20%)
%    - Conference proceedings (10-20%)
%    - Internet sources (5-15%, hanya yang kredibel)
%    - Grey literature (reports, theses) - minimal 5%
%
% 3. FORMAT MENGGUNAKAN BIBTEX:
%    
%    a) Buat file references.bib:
%       @article{Author2020,
%         title={Title of Paper},
%         author={Author, A. and Others, B.},
%         journal={Journal Name},
%         volume={30},
%         pages={100-115},
%         year={2020},
%         doi={10.xxxx/xxxxx}
%       }
%
%    b) Cite dalam dokumen: \citet{Author2020} atau \citep{Author2020}
%
% 4. TOOLS UNTUK MANAJEMEN REFERENSI:
%    - Mendeley (mendeley.com) - generate .bib file
%    - Zotero (zotero.org) - open source
%    - EndNote - commercial
%    - RefWorks - web-based
%
% 5. GENERATOR BIB FORMAT:
%    - Google Scholar: export ke BibTeX
%    - Semantic Scholar: export BibTeX
%    - CrossRef: provide DOI, convert ke BibTeX
%
% 6. TIPS PENTING:
%    - Gunakan DOI jika tersedia
%    - Pastikan spelling penulis dan tahun AKURAT
%    - Urutan publikasi: Chronological atau Alphabetical
%    - Dalam disiplin ilmu tertentu, format mungkin berbeda
%      (tanyakan pembimbing tentang format yang diharapkan)
% ====================================================================
